\documentclass[11pt]{article}
\usepackage[margin=.85in]{geometry}
\usepackage[T1]{fontenc}
\usepackage{lmodern,amsmath,amssymb,amsthm,booktabs,microtype,xcolor,hyperref,fancyhdr}
\hypersetup{colorlinks=true,linkcolor=blue!40!black,urlcolor=blue!45!black,
pdftitle={Principal centering and the remaining signed correlation},pdfauthor={Jude Gomila}}
\pagestyle{fancy}\fancyhf{}\fancyhead[L]{\small ZETA EXPLORER / PRINCIPAL SECTOR}
\fancyhead[R]{\small October 8, 2026}\fancyfoot[C]{\thepage}
\setlength{\headheight}{14pt}
\newtheorem{theorem}{Theorem}\newtheorem{lemma}{Lemma}\newtheorem{proposition}{Proposition}
\newcommand{\dd}{\,d}\newcommand{\E}{\mathcal E}\newcommand{\C}{\mathcal C}
\newcommand{\R}{\mathcal R}\newcommand{\D}{\mathcal D}
\title{Principal centering and the remaining signed correlation\\[4pt]
\large A continuation of the quasi-RH moment-transfer audit}
\author{Research continuation for Jude Gomila}\date{October 8, 2026}
\begin{document}\maketitle
\begin{abstract}
We carry out the proposed principal-sector reduction for the existing
Vaughan packet. The diagonal is bounded explicitly by a degree-five
polynomial in the logarithmic scale, and all shifts of integer size at
most $H$ cost at most $2H$ times that bound. We compute the full principal
Laurent part, including the derivative of the truncated M\"obius series,
and give an exact centered Stieltjes representation. The packet already
annihilates this principal density. The rational-to-Eisenstein lift
retains a product cutoff and leaves only one unweighted variable, of
length at most $e^{T/25+3/2}$. These facts identify why the cited
equal-product rectangle estimate does not directly apply. No positive
exponent saving is proved. We state the remaining signed large-shift
estimate precisely and verify the finite identities with exact arithmetic
and 512-bit Arb calculations. This is a working research note, not an
independently refereed or formally verified theorem, and no novelty is claimed.
\end{abstract}

\section{The fixed packet and the uniform target}
Use the real even packet $k$ from the v9 note and the preceding transfer
audit \cite{v9,transfer}. It satisfies $\operatorname{supp}k\subset[-1/2,1/2]$,
$|k|\le1$, $k(0)=1$, and
\begin{equation}\label{K}
 \begin{aligned}
 K(w)&=\int k(t)e^{wt}\dd t
 =\frac{(1/4-w^2)^4}{\kappa}
       \left(\frac{\sinh(w/48)}{w/48}\right)^{24},\\
 \kappa&=\frac{11431397901559182284554757008547201}{2872446304320552960000}.
 \end{aligned}
\end{equation}
In particular $K^{(j)}(1/2)=0$ for $0\le j\le3$, while
$K(w)\ne0$ for $0<\Re w<1/2$. All sums below include every prime power.

For each real $T\ge1$, hold $U=V=\lfloor e^{12T/25}\rfloor$ fixed
throughout $L\in[T,T+1]$. Define
\begin{align}
 a_U(r)&=\sum_{d\mid r,\ d>U}\mu(d), &
 d_T(n)&=\sum_{rb=n,\ b>V}a_U(r)\Lambda(b),\label{d}\\
 D_T(L)&=\sum_n\frac{d_T(n)}{\sqrt n}k(\log n-L),&
 \E_{\rm II}(T)&=\int_T^{T+1}|D_T(L)|^2\dd L.\label{energy}
\end{align}
The condition $r>U$ is automatic if $a_U(r)\ne0$.
The inherited norm error between this packet and the prime packet is
$6\cdot10^{15}(T+2)e^{-T/50}$. Consequently their nonnegative logarithmic
energy exponents agree. The preceding note derives from the externally
posted $7/8$ assertion the bound $\E_{\rm II}(T)=O(e^{3T/4})$.
We continue to treat that assertion as an external input; its complete
Lean project has not been rebuilt here.

\section{The diagonal and nearby integer shifts}
Put $J_T=[e^{T-1/2},e^{T+3/2}]\cap\mathbb N$, set
$z_n=d_T(n)/\sqrt n$ on $J_T$ and zero elsewhere, and write
\[
 G_T(m,n)=\int_T^{T+1}k(\log m-L)k(\log n-L)\dd L.
\]
This real Gram matrix is positive semidefinite,
$|G_T(m,n)|\le c_k:=\|k\|_2^2\le1$, and its entries vanish when
$|\log(m/n)|\ge1$. Individual off-diagonal entries need not be positive.
For an integer $H\ge0$, let
\begin{align*}
 A(T)&=\sum_n z_n^2G_T(n,n),\\
 \C_{\le H}(T)&=2\sum_{\substack{m<n\\n-m\le H}}z_mz_nG_T(m,n),&
 \C_{>H}(T)&=2\sum_{\substack{m<n\\n-m>H}}z_mz_nG_T(m,n).
\end{align*}
Exactly, $\E_{\rm II}=A+\C_{\le H}+\C_{>H}$. The grouping by the
\emph{integer product} $n$ precedes this decomposition: different
factorizations of one $n$ are not independent diagonal entries.

\begin{theorem}[Explicit diagonal and short-shift bounds]\label{band}
For every real $T\ge1$, every fixed positive integer pair $U,V$, and
every integer $H\ge0$, define
\[
 Q(T)=e^2(T+3/2)^2(T+5/2)^3.
\]
Then
\begin{align}
 \sum_{n\in J_T}\frac{|d_{U,V}(n)|^2}{n}&\le Q(T),\label{mass}\\
 0\le A(T)&\le c_k Q(T),\label{diag}\\
 2\sum_{\substack{m<n\\n-m\le H}}|z_mz_nG_T(m,n)|
 &\le 2Hc_k Q(T),\label{near}\\
 \C_{>H}(T)&\ge-(2H+1)c_k Q(T).\label{lower}
\end{align}
In particular these bounds remain uniform for the moving window cutoffs.
\end{theorem}
\begin{proof}
Since $|a_U(r)|\le\tau(r)\le\tau(n)$ for $r\mid n$ and
$\sum_{b\mid n}\Lambda(b)=\log n$, we have
$|d_{U,V}(n)|\le\tau(n)\log n$. At a prime power $p^j$,
$(j+1)^2\le\binom{j+3}{3}$; multiplicativity gives
$\tau(n)^2\le\tau_4(n)$. Counting ordered quadruples yields
\[
 \sum_{n\le x}\tau_4(n)
 \le x\left(\sum_{a\le x}\frac1a\right)^3
 \le x(1+\log x)^3.
\]
Take $x=e^{T+3/2}$ and bound $1/n\le e^{-T+1/2}$ on $J_T$ to
obtain \eqref{mass}. The Gram bound proves \eqref{diag}.
For each shift $1\le h\le H$, extending $z$ by zero and using
$2|z_nz_{n+h}|\le z_n^2+z_{n+h}^2$ gives
$\sum_n|z_nz_{n+h}|\le\sum_nz_n^2$.
Sum this inequality over $h$ to obtain \eqref{near}.
Finally $\E_{\rm II}\ge0$ implies
$\C_{>H}\ge-A-\C_{\le H}\ge-(2H+1)c_kQ$.
\end{proof}
This controls the negative part of the far correlation. The remaining
task is its \emph{upper} bound. Positivity of the full energy supplies
no such bound. Cauchy--Schwarz over every pair gives only an
$O(e^T(1+T)^5)$ envelope, not an improvement on the imported result.

\section{The exact principal density}
For this section fix $U,V$ before taking transforms. Write
\[
 M(s)=\sum_{d\le U}\mu(d)d^{-s},\quad
 L_V(s)=\sum_{b\le V}\Lambda(b)b^{-s},\quad
 m_0=M(1),\quad m_1=M'(1),\quad \ell_0=L_V(1).
\]
The exact identity from the transfer audit is
\begin{equation}\label{Didentity}
 \D_{U,V}(s)=\sum_n d_{U,V}(n)n^{-s}
 =-\frac{\zeta'}\zeta(s)-L_V(s)+M(s)\zeta'(s)
       +\zeta(s)M(s)L_V(s).
\end{equation}

\begin{proposition}[Full principal part and centered response]\label{center}
With $B=1-m_1+m_0\ell_0$, the principal part at $s=1$ is
\begin{equation}\label{principal}
 -\frac{m_0}{(s-1)^2}+\frac{B}{s-1}.
\end{equation}
Define
\begin{align*}
 p_{U,V}(x)&=-m_0\log x+B,\\
 P_{U,V}(x)&=\int_1^xp_{U,V}(t)\dd t
 =x[-m_0\log x+m_0+B]-m_0-B,\\
 \R_{U,V}(x)&=\sum_{n\le x}d_{U,V}(n)-P_{U,V}(x).
\end{align*}
For every $L\ge1$, the following identities are exact:
\begin{align}
 \int_1^\infty p_{U,V}(x)x^{-1/2}k(\log x-L)\dd x&=0,\label{densityzero}\\
 D_{U,V}(L)&=\int_1^\infty\R_{U,V}(x)x^{-3/2}
       [\tfrac12 k(\log x-L)-k'(\log x-L)]\dd x.\label{Stieltjes}
\end{align}
\end{proposition}
\begin{proof}
At $s=1+z$, $\zeta(s)=z^{-1}+O(1)$ and
$\zeta'(s)=-z^{-2}+O(1)$. Multiplying $M(1+z)=m_0+m_1z+O(z^2)$
in \eqref{Didentity} gives \eqref{principal}. In particular the term
$-m_1$ must not be omitted. The Mellin integral of $p$ on $[1,\infty)$
has exactly this principal part.

Set $t=\log x-L$. The support lies above $x=1$ for $L\ge1$,
and the density integral equals
\[
 e^{L/2}\bigl[(B-m_0L)K(1/2)-m_0K'(1/2)\bigr]=0.
\]
For $f_L(x)=x^{-1/2}k(\log x-L)$,
$f_L'(x)=x^{-3/2}[k'(\log x-L)-k(\log x-L)/2]$.
Now $D=\int f_L\dd\R+\int f_Lp\dd x=-\int\R f_L'\dd x$.
The compact packet vanishes at both endpoints, so the integration-by-parts
boundary term is zero. On a truncated subinterval it is generally
$[\R f_L]_a^b$ and cannot be discarded.
\end{proof}
These identities also apply separately on each window with its frozen
$U(T),V(T)$. They do not justify substituting moving cutoffs into a
single Dirichlet-series continuation argument.

The principal density is therefore already removed by the existing
packet. At a nontrivial zero $\rho$ of multiplicity $m_\rho$, subtracting
\eqref{principal} leaves the residue $-m_\rho$ unchanged. Multiplying
by $K(\rho-1/2)$ does not remove a zero with $\Re\rho>1/2$.
Thus this correct centering supplies an exact residual, not an exponent saving.

\section{The coefficient lift and the short remaining plain variable}
Let $\chi=\chi_{-3}$ and let $\nu(n)$ be the sum of ideal M\"obius
coefficients over ideals of norm $n$ in $\mathbb Q(\sqrt{-3})$.
The established identities are $\nu=\mu*(\mu\chi)$ and
$\mu=\nu*\chi$. Expanding the divisor in \eqref{d} gives
\begin{equation}\label{lift}
 d_{U,V}(n)=\sum_{dbq=n\atop d>U,\ b>V}\mu(d)\Lambda(b)
 =\sum_{acbq=n\atop ac>U,\ b>V}\nu(a)\chi(c)\Lambda(b).
\end{equation}
The cutoff is $ac>U$, not $a>U$. Inserting this expression twice into
the Gram formula retains the two product cutoffs and every mixed term.
For clarity, with $n=acbq$ and $n'=a'c'b'q'$, the weight is
\[
 \frac{\nu(a)\nu(a')\chi(c)\chi(c')\Lambda(b)\Lambda(b')}
 {\sqrt{nn'}}G_T(n,n'),
 \qquad ac,a'c'>U,\quad b,b'>V.
\]
The far-shift restriction is additionally $|n-n'|>H$; no character
average is present in this exact identity.

For the balanced cutoffs, every nonzero term on the window satisfies
\begin{equation}\label{short}
 q=\frac{n}{db}< e^{T/25+3/2},
 \qquad d,b\ge U+1>e^{12T/25}.
\end{equation}
The norm lift splits $d=ac$ but does not create another unweighted
long variable. If
$G(y)=\sum_{q\ge1}q^{-1/2}k(\log q-y)$, the exact alternative form is
\[
 D_T(L)=\sum_{d>U,\ b>V}\frac{\mu(d)\Lambda(b)}{\sqrt{db}}
                  G(L-\log(db)).
\]
The ordinary continuum moment vanishes, but the actual discrete $q=1$
term remains. The outer product can reach $db\asymp e^L$, so this
formula cannot be bounded by assuming a large positive argument of $G$
uniformly over its outer sum.

\section{Attempted use of the matched rectangle estimate}
The pinned external paper's \texttt{lem:centered-lattice-cancellation}
concerns plain lattice sums with one common finite-order character,
one common puncture mask, fixed profiles, and one common norm power.
Their leading terms have the form
$cX^{1+it}I_i(t)$. Products at scales satisfying
$X_1X_2=Y_1Y_2$ have identical leading terms. If each remaining length
is at least $Z^r$, the stated normalized centered bound is
\[
 O\bigl(Z^{\epsilon_1}(1+|t|)^{2J}
                (X_1X_2)^{1/2}Z^{-r}\bigr).
\]
The preceding \texttt{lem:centered-coefficient-invariant} is what
preserves those common data through the source's transforms.
This is not an assertion about every centered arithmetic expression.

Our attempt reaches the following exact comparison:
\begin{center}\small
\begin{tabular}{p{.27\linewidth}p{.65\linewidth}}
\toprule
Required source structure & What the present reduction supplies\\
\midrule
Two plain variables with matching character and mask &
$\nu(a)\chi(c)\Lambda(b)$ and one unweighted variable $q$; the latter
has only the upper length bound \eqref{short}.\\[3pt]
Equal-product rectangle difference &
A hyperbolic cutoff $ac>U$, $b>V$, followed by a product-difference
restriction $|acbq-a'c'b'q'|>H$.\\[3pt]
Coefficient invariant under both transforms &
No such transformed identity has been constructed for this quadratic form.\\[3pt]
Uniform reconstruction and remainder &
The principal density is removed exactly, but no small reconstruction
remainder or uniform exceptional-row energy bound follows.\\
\bottomrule
\end{tabular}
\end{center}
The source's \texttt{lem:plain} still excludes the principal inducing
row (and a larger finite family when live prime slots are present).
The source separately handles a principal numerator in its own Mellin
signal. That residue calculation is not a license to insert our
principal arithmetic sequence into the excluded-row estimate.

Trying to sum absolute pair contributions gives the $e^T$ envelope
above. Applying the imported zero bound to the original packet and
transferring its norm gives the existing $e^{3T/4}$ bound. Neither
method yields a positive $\delta$. A new exact transformed identity
and a compatible signed bound are needed to use the matched rectangle
mechanism here. This is a limitation of this attempted transfer, not
an impossibility theorem for arithmetic cancellation.

\paragraph{A control against centering alone.}
For a synthetic density $x^{\beta-1}$ on $(0,\infty)$, filtering gives
\[
 \int_0^\infty x^{\beta-1}x^{-1/2}k(\log x-L)\dd x
 =K(\beta-1/2)e^{(\beta-1/2)L}.
\]
At $\beta=55/64<7/8$, its energy grows with exponent $23/32$.
Adding or subtracting any density $A\log x+B$ makes no difference.
It therefore violates the illustrative improved exponent $11/16$
despite satisfying the same principal-moment cancellation. This is
a synthetic control, not a counterexample involving zeta or its
M\"obius coefficients.

\section{A precise next theorem, and what it would imply}
Let $H(T)=\lfloor e^{T/2}\rfloor$ and choose a fixed
$0<\delta<1/4$. Put $\alpha=3/4-\delta>1/2$.
The concrete remaining assertion is
\begin{equation}\label{target}
 \boxed{\C_{>H(T)}(T)\le M_\delta e^{\alpha T}
 \quad\text{for all real }T\ge T_\delta,}
\end{equation}
for constants independent of $T$ and its prescribed cutoffs.
This is a one-sided signed estimate; it does not request a bound on
the sum of absolute values or on each dyadic shift band separately.

\begin{proposition}[Reduction to the large shifts]
For the stated $\alpha$, \eqref{target} holds if and only if
$\E_{\rm II}(T)=O(e^{\alpha T})$.
\end{proposition}
\begin{proof}
Theorem~\ref{band} gives
$|A+\C_{\le H}|\le(2H+1)c_kQ(T)
=O(e^{T/2}(1+T)^5)=O(e^{\alpha T})$.
Subtract this term from, or add it to,
$\E_{\rm II}=A+\C_{\le H}+\C_{>H}$.
\end{proof}
Thus the target is an exact localization of the difficulty, not an
independently easier theorem already established. The v8/v9 analytic
growth identity, retained in the preceding audit, is
\[
 \limsup_{T\to\infty}\frac{\log(1+\E_{\rm II}(T))}{T}
       =2\sup_\rho(\Re\rho-1/2).
\]
Using that identity, \eqref{target} would imply
$\Re\rho\le7/8-\delta/2$ for every nontrivial zero.
In the shifted-$\xi$ bridge interpretation it would lower the sufficient
width from $3/8$ to $3/8-\delta/2$.
The inherited growth and bridge arguments still require independent review.

For example $\delta=1/16$ would mean an energy exponent $11/16$ and
a zero boundary $27/32$. These numbers illustrate a possible target;
they are not a proved or predicted saving. The fixed choice
$H=e^{T/2}$ is for a first improvement only. To iterate toward RH one
would also have to lower the near-shift exponent and control all
constants and hypotheses uniformly. No such iteration is established.

\section{Verification and research status}
\texttt{audit.py} uses only exact integers and rational numbers. It checks
16,384 formal-log coefficients in four cutoff configurations, including
the full norm lift and its cutoff; local divisor bounds; twelve generic
principal Laurent models; and twelve rational Gram decompositions.
Replacing $ac>U$ by $a>U$, dropping $M'(1)$, and mismatching the two
rectangle masks are deliberately rejected.

\texttt{interval\_audit.py} uses \texttt{python-flint==0.8.0} at 512 bits.
It integrates the actual packet by exact piecewise polynomial or
exponential-polynomial antiderivatives evaluated in Arb, not heuristic
quadrature tolerances. It checks the four vanishing moments, a nonzero
off-axis response, the signed energy decomposition, and the centered
Stieltjes identity at actual arithmetic windows. The omitted-$k'$
control is rejected. All coefficients, including repeated prime powers,
are retained. The detailed enclosures and hashes are in the receipts.

The following rounded values illustrate why neither sign of the far
correlation can simply be assumed; the receipts contain rigorous balls.
\begin{center}\small
\begin{tabular}{rrrrr}
\toprule
$T$ & $U$ & $H$ & $\E_{\rm II}(T)$ & $\C_{>H}(T)$\\
\midrule
2 & 2 & 2 & $0.057468669$ & $-0.002863840$\\
3 & 4 & 4 & $0.087225379$ & $+0.010084453$\\
4 & 6 & 7 & $0.093690267$ & $-0.007399596$\\
\bottomrule
\end{tabular}
\end{center}
These finite samples do not certify a bound at unbounded scales.
The analytic estimates in Theorem~\ref{band} and Proposition~\ref{center}
are established by their displayed proofs, not inferred from the tests.

\textbf{Outcome.} We have proved an explicit diagonal/near-shift estimate,
completed the principal-density subtraction, retained the exact lift,
and isolated a signed large-shift target. The attempted external
rectangle transfer does not close: its coefficient invariant and
uniform remainder are missing. \textbf{No positive $\delta$, improved
zero-free region, or RH proof has been obtained.}

\begin{thebibliography}{9}\small
\bibitem{transfer} \emph{Testing the quasi-RH moment transfer}, research
continuation for Jude Gomila, October 6, 2026. Archived in
\texttt{research-inputs/rh-moment-transfer-2026-10-06/}, Zeta Explorer PR 11.
\bibitem{v9} \emph{Vaughan reduction of the filtered prime energy},
October 1, 2026, user bundle \texttt{rh\_vaughan\_v9.zip}, including its
retained v8 growth argument. Exact input hashes are in the source index.
\bibitem{source} OpenAI, \emph{The Quasi-Riemann Hypothesis: A Zero-Free
Half-Plane $\Re(s)>7/8$}, posted September 30, 2026. Source pinned to
commit \texttt{adc7f1241b42e322a6451854ab7e4b4c146bf78a}.
Labels used: \texttt{lem:plain}, \texttt{old-eq:2.18a},
\texttt{lem:centered-coefficient-invariant}, and
\texttt{lem:centered-lattice-cancellation}.
\url{https://github.com/openai/math/blob/main/lean/docs/003.md}.
\end{thebibliography}
\end{document}
